\chapter{Công nghệ sử dụng}
\label{ch:tech}

Chương này giới thiệu các công nghệ, nền tảng và công cụ chính được sử dụng để xây dựng hệ
thống \brandName{}, kèm theo lý do lựa chọn của từng công nghệ trong ngữ cảnh của đề tài.

\section{Next.js, React và TypeScript}

\subsection{React}

React là một thư viện JavaScript mã nguồn mở do Meta phát triển, dùng để xây dựng giao diện
người dùng theo mô hình thành phần (component). Thay vì thao tác trực tiếp với DOM, React mô
tả giao diện một cách khai báo và sử dụng cơ chế DOM ảo (Virtual DOM) để cập nhật giao diện
một cách hiệu quả khi dữ liệu thay đổi. Cách tiếp cận theo thành phần giúp mã nguồn dễ tái sử
dụng, dễ kiểm thử và dễ bảo trì \cite{react,reactjs2018}.

\subsection{Next.js 14 với App Router}

Next.js là một framework xây dựng trên nền React, cung cấp một bộ công cụ hoàn chỉnh cho việc
phát triển ứng dụng web full-stack. Đề tài sử dụng \textbf{Next.js phiên bản 14} với mô hình
\textit{App Router} --- mô hình định tuyến dựa trên cấu trúc thư mục, hỗ trợ nhiều cơ chế kết
xuất trang khác nhau \cite{nextjs}:

\begin{itemize}
  \item \textbf{Server-Side Rendering (SSR):} trang được kết xuất phía máy chủ ở mỗi yêu cầu,
  phù hợp với nội dung động và thân thiện với SEO.
  \item \textbf{Static Site Generation (SSG):} trang được dựng sẵn tại thời điểm build, cho tốc độ
  tải rất nhanh.
  \item \textbf{Incremental Static Regeneration (ISR):} kết hợp ưu điểm của hai cơ chế trên, cho
  phép tái sinh trang tĩnh theo chu kỳ hoặc theo nhãn (tag) dữ liệu.
\end{itemize}

Điểm nổi bật của App Router là \textbf{React Server Components (RSC)} và \textbf{Server
Actions}. Với RSC, phần lớn các thành phần được kết xuất hoàn toàn trên máy chủ và không gửi
mã JavaScript xuống trình duyệt, giúp giảm đáng kể kích thước gói tải về và cải thiện hiệu năng.
Server Actions cho phép định nghĩa các hàm xử lý nghiệp vụ chạy trên máy chủ và gọi trực tiếp
từ thành phần giao diện như một hàm thông thường, loại bỏ nhu cầu xây dựng nhiều endpoint API
thủ công cho các thao tác ghi dữ liệu. Đây là lý do chính khiến Next.js được lựa chọn: nó cho phép
xây dựng một ứng dụng full-stack thống nhất, vừa tối ưu SEO và hiệu năng --- hai yếu tố then chốt
của một website thương mại điện tử --- vừa rút ngắn thời gian phát triển.

\subsection{TypeScript}

Toàn bộ mã nguồn của hệ thống được viết bằng \textbf{TypeScript} ở chế độ kiểm tra kiểu nghiêm
ngặt (strict). TypeScript bổ sung hệ thống kiểu tĩnh cho JavaScript, giúp phát hiện sớm nhiều lỗi
ngay trong quá trình phát triển, tăng khả năng tự gợi ý của trình soạn thảo và giúp mã nguồn dễ
bảo trì hơn khi dự án mở rộng \cite{typescript}.

\section{Prisma, PostgreSQL và Supabase}

\subsection{PostgreSQL}

PostgreSQL là một hệ quản trị cơ sở dữ liệu quan hệ - đối tượng mã nguồn mở, nổi bật về tính ổn
định, tuân thủ chuẩn SQL và khả năng mở rộng. Hệ thống \brandName{} khai thác một số tính năng
nâng cao của PostgreSQL \cite{postgresql}:

\begin{itemize}
  \item \textbf{Tìm kiếm toàn văn (Full-Text Search)} kết hợp tiện ích mở rộng \texttt{unaccent}
  để loại bỏ dấu tiếng Việt và \texttt{pg\_trgm} để so khớp gần đúng, phục vụ chức năng tìm kiếm
  sản phẩm.
  \item \textbf{Row-Level Security (RLS)} để thiết lập các chính sách truy cập dữ liệu ở mức dòng,
  tăng cường an toàn dữ liệu.
  \item Hỗ trợ kiểu dữ liệu phong phú như \texttt{uuid}, \texttt{jsonb}, mảng (array) và kiểu liệt
  kê (enum).
\end{itemize}

\subsection{Supabase}

Supabase là một nền tảng Backend-as-a-Service xây dựng quanh PostgreSQL, cung cấp sẵn cơ sở
dữ liệu được quản lý, dịch vụ xác thực (Authentication), lưu trữ tệp (Storage) và các API. Trong
đề tài, Supabase đảm nhiệm vai trò hạ tầng cơ sở dữ liệu và xác thực người dùng, đồng thời cung
cấp dịch vụ lưu trữ hình ảnh sản phẩm và ảnh đánh giá \cite{supabase}. Việc sử dụng Supabase giúp
tập trung nguồn lực vào phát triển nghiệp vụ thay vì tự xây dựng và vận hành hạ tầng.

\subsection{Prisma ORM}

Prisma là một ORM (Object-Relational Mapping) hiện đại cho hệ sinh thái Node.js/TypeScript.
Lược đồ cơ sở dữ liệu được mô tả tập trung trong tệp \texttt{schema.prisma}, từ đó Prisma sinh ra
một thư viện truy vấn an toàn kiểu (type-safe client) \cite{prisma}. Ưu điểm khi sử dụng Prisma:

\begin{itemize}
  \item Truy vấn an toàn kiểu, hạn chế lỗi và tận dụng tối đa khả năng gợi ý của TypeScript.
  \item Quản lý phiên bản lược đồ cơ sở dữ liệu thông qua cơ chế migration.
  \item Hỗ trợ giao dịch (transaction) với các mức cô lập khác nhau, đặc biệt quan trọng đối với
  các nghiệp vụ nhạy cảm như đặt hàng và trừ tồn kho.
\end{itemize}

\section{Tích hợp dịch vụ bên ngoài}

\subsection{Cổng thanh toán PayOS}

PayOS là một cổng thanh toán trực tuyến phổ biến tại Việt Nam, hỗ trợ thanh toán qua mã QR và
chuyển khoản ngân hàng. Hệ thống tích hợp PayOS để tạo liên kết thanh toán cho đơn hàng, nhận
kết quả thanh toán thông qua webhook có xác thực chữ ký HMAC, và tự động hủy các đơn hàng
chưa thanh toán quá hạn nhằm hoàn trả tồn kho \cite{payos}.

\subsection{Resend và Upstash}

\textbf{Resend} là dịch vụ gửi email giao dịch qua API, được sử dụng để gửi các email thông báo
như xác nhận đặt hàng, cập nhật trạng thái đơn hàng và xác nhận thanh toán. \textbf{Upstash
Redis} là dịch vụ Redis serverless, được sử dụng để giới hạn tần suất truy cập (rate limiting) cho
các thao tác nhạy cảm như đăng nhập và đăng ký, góp phần chống tấn công dò mật khẩu. Cả hai
dịch vụ đều được tích hợp theo cơ chế ``env-gated'': nếu không cấu hình khóa API, tính năng tương
ứng sẽ tự động bỏ qua mà không ảnh hưởng đến hoạt động chung của hệ thống.

\section{Thư viện giao diện và công cụ hỗ trợ}

\begin{itemize}
  \item \textbf{Tailwind CSS:} framework CSS theo hướng utility-first, giúp xây dựng giao diện
  nhất quán và nhanh chóng \cite{tailwind}.
  \item \textbf{Zustand:} thư viện quản lý trạng thái phía client nhẹ, dùng cho giỏ hàng, danh sách
  yêu thích và so sánh sản phẩm.
  \item \textbf{React Hook Form và Zod:} quản lý biểu mẫu và kiểm tra dữ liệu đầu vào; Zod được
  dùng để định nghĩa lược đồ kiểm tra dùng chung cho cả phía client và phía máy chủ.
  \item \textbf{SheetJS (xlsx):} xuất báo cáo doanh thu ra tệp Excel.
  \item \textbf{Recharts} và \textbf{dnd-kit:} lần lượt phục vụ vẽ biểu đồ báo cáo và thao tác kéo
  thả sắp xếp banner trong trang quản trị.
\end{itemize}

\section{Nền tảng triển khai}

Hệ thống được triển khai trên \textbf{Vercel} --- nền tảng được tối ưu cho ứng dụng Next.js, hỗ trợ
triển khai liên tục (CI/CD) từ kho mã nguồn, chạy các hàm serverless và tác vụ định kỳ (cron
job). Trong đề tài, một tác vụ cron được cấu hình để định kỳ quét và hủy các đơn hàng PayOS chưa
thanh toán quá hạn \cite{vercel}.
